\section{Quantum records and the emergence bridge}

Quantum theory supplies the sharpest positive precedents for information
leaving its first bearer, and the sharpest warning against treating that
precedent as a derivation of quantum theory.  Once a Hilbert space, a tensor or
algebraic decomposition, a state, and an intervention class are supplied,
standard quantum mechanics already explains how interactions distribute
correlations, how some observables become stable under environmental
monitoring, and how collective excitations acquire particle-like roles.  The
IOTA hypothesis begins one level earlier.  It must explain how its local
informational distinctions could acquire the composition, phase, probability,
and transformation structure presupposed by those results.

Table~\ref{tab:emergence-ladder} makes this separation explicit.  The first
rung, and precursor aspects of the second, have bounded IOTA evidence; no
project result yet establishes quantum coherence.  The current Dirac
construction occupies a higher, supplied-quantum rung.  The remaining arrows
are research hypotheses, not consequences of similarity in vocabulary.

\begin{table}[t]
\centering
\small
\caption{The emergence ladder and the evidential status of each bridge.}
\label{tab:emergence-ladder}
\begin{tabularx}{\linewidth}{@{}>{\raggedright\arraybackslash}p{0.16\linewidth}
  >{\raggedright\arraybackslash}X >{\raggedright\arraybackslash}X@{}}
\toprule
Rung & Present support & Missing discriminating result \\
\midrule
Local interaction $\rightarrow$ transported constraint
& Project construction: selected connectivity roles are re-instantiated in
disjoint successors, with causal intervention gates.
& Replication across rules, schedules, topologies, and physically equivalent
representations. \\
Transport $\rightarrow$ persistent coherence
& Project lineages sometimes outlive original addresses, a pre-quantum
persistence result; quantum and biological records are independent precedents.
& A minimal compositional invariant whose persistence is not an artifact of a
chosen tracking rule. \\
Coherence $\rightarrow$ quantum structure
& No base-IOTA derivation.  Process tensors, decoherence, reference frames, and
the Dirac project all assume quantum theory.
& Derive state composition, reversible/interventional dynamics, interference,
and probability calibration without importing amplitudes or the Born rule. \\
Quantum structure $\rightarrow$ quasiparticle
& Established local quantum models and experiments realize fractionalized and
topological excitations.
& Create and manipulate an IOTA-derived excitation with declared dispersion,
statistics, symmetry, conservation, and scattering. \\
Quasiparticle $\rightarrow$ fields, mass, and energy
& One-dimensional quantum cellular automata recover bounded free-field roles;
the project recovers a bounded one-particle Dirac role.
& Interacting multiparticle limits, no-go closure, and derived rather than
scheduled mass- and energy-like invariants. \\
\bottomrule
\end{tabularx}
\end{table}

\subsection{Multi-time quantum memory is a native future-law formalism}

\paragraph{Established field result.}
The process-tensor framework represents the statistics of a quantum process
under a sequence of possible interventions.  It separates the experimenter's
operations from the multi-time object that maps those operations to outcome
probabilities; quantum Markov order and memory classifications then depend on
the chosen instruments and system--environment split
\citep{PollockEtAl2018ProcessTensor,PollockEtAl2018Markov,
TarantoEtAl2019MarkovOrder,TarantoEtAl2024ClassicalMemory}.  Recent process
tomography demonstrates that such objects can be reconstructed in controlled
devices, subject to finite-dimensional, sampling, and control assumptions
\citep{GiarmatziEtAl2025Tomography}.  This is close to
Eq.~\eqref{eq:future-family}: the future object is not one unconditioned state
but a family of responses to allowed operations.

\paragraph{Synthesis inference.}
A multi-time process is a natural quantum benchmark for an ICP carrier claim.
Two histories can leave the same measured reduced system state yet differ in
environmental correlations and therefore in their later intervention-response
families.  The record has not escaped physics; it resides in an enlarged joint
state or process.  Consequently, state enlargement is both a positive method
for locating memory and the strongest null against a new historical variable.

\paragraph{ICP hypothesis.}
The stronger claim is that a compact, physically identifiable role can be
formed, edited, and transferred while preserving the relevant process family.
That claim would require histories matched on the best available present state,
randomized preservation or rewriting of the proposed carrier, resource-matched
shams, and held-out multi-time interventions.  If the augmented standard
process tensor explains every effect without a new invariant or compression,
ICP remains a useful organization of ordinary quantum memory rather than an
additional quantum law.

\subsection{Environmental records: a real mechanism with an objectivity ladder}

\paragraph{Established field result.}
Decoherence disperses phase information into correlations with an environment,
and environment-induced superselection identifies observables comparatively
stable under the specified interaction.  Quantum Darwinism adds a distinct
condition: many disjoint fragments may carry information about a common
preferred observable.  Environment-as-witness constructions and a broad
finite-dimensional channel theorem show how fragment channels can become
approximately measure-and-prepare channels for one common POVM
\citep{OllivierPoulinZurek2004Witness,
BrandaoEtAl2015Darwinism}.  Controlled photonic and spin experiments
have observed redundant imprints under engineered or well-characterized
conditions \citep{CiampiniEtAl2018DarwinismExperiment,ChenEtAl2019DarwinismExperiment,
UndenEtAl2019Classicality}.

This is direct evidence that an interaction-generated distinction can be
embodied in secondary degrees of freedom and read away from the source.  It is
not a demonstration of information without a carrier.  Nor do the familiar
levels collapse into one another:
\begin{equation}
 \text{decoherence}\not\equiv\text{pointer selection}
 \not\equiv\text{redundant encoding}
 \not\equiv\text{accessible consensus}
 \not\equiv\text{objective broadcast record}.
 \label{eq:objectivity-ladder}
\end{equation}
Spectrum broadcast structure supplies a stringent state-level form in which
multiple fragments have distinguishable conditional supports.  Strong quantum
Darwinism additionally requires locally accessible classical information, and
the multipartite-independence condition must be stronger than pairwise
conditional independence \citep{HorodeckiEtAl2015Objectivity,
LeOlayaCastro2018Objectivity,FellerEtAl2021CommentDarwinism}.  Mutual-information
plateaus, redundancy, and observer consensus can therefore disagree
\citep{ChisholmEtAl2023RedundancyConsensus}.  Classification results further show that
record proliferation requires dynamical structure; arbitrary coupling does not
automatically make an objective record \citep{DoucetDeffner2024DarwinismHamiltonians}.

\paragraph{Synthesis inference.}
Quantum Darwinism supplies the clearest independent model of persistence beyond
an original bearer: correlations generated at a source can remain available in
other physical carriers.  It does not yet supply the ICP succession criterion.
Most experiments ask what can be inferred from a fragment now, not whether
erasing, swapping, or replacing that fragment later changes a registered
future process while all ordinary resources are controlled.

\paragraph{Factorization boundary.}
The words ``system,'' ``environment,'' and ``fragment'' presuppose a
decomposition.  The same global state can display redundant imprints of
noncommuting observables under different, non-comparable environment
partitions \citep{Fu2021RedundantImprints}; a fixed Hamiltonian can also admit
multiple tensor-product structures with robust pointer states that need not
track familiar locality \citep{AdilEtAl2026ClassicalSubsystems}.  This does not
undermine an operationally fixed laboratory partition.  It blocks the stronger
inference from a redundant record to a unique, observer-independent carrier or
classical world.  An ICP quantum claim must therefore state the factorization
or observable algebra, accessible fragment operations, reference resources,
and the transformations under which carrier identity is expected to survive.

\subsection{Relational facts and quantum reference frames are not memories}

Relational quantum mechanics (RQM) and quantum reference-frame (QRF) theory
share an insistence that descriptions are relational, but they answer different
questions.  RQM is an interpretation in which variable values are facts
relative to interactions; later work distinguishes approximately stable facts
from generic relative facts \citep{Rovelli1996RelationalQuantum,
DiBiagioRovelli2021StableFacts}.  Whether its commitments permit consistent
cross-perspective fact assignments remains actively disputed
\citep{Pienaar2021Quandaries,DiBiagioRovelli2022Reply}.  That dispute is
interpretive and supplies no independent experimental confirmation of IOTA.

QRF formalisms instead transform states, observables, dynamics, and
measurements relative to physical quantum systems.  Superposition and
entanglement can be frame-relative while consistently transformed operational
probabilities remain invariant \citep{GiacominiEtAl2019QuantumFrames}.  Operational
equivalence and reversibility depend on the framework, localization, symmetry
constraints, and access to global charges
\citep{CaretteEtAl2025OperationalFrames,DoatVanrietvelde2026Charge}.

\paragraph{Synthesis inference.}
This literature requires ICP to distinguish a changed description from a
changed future.  A frame transformation may relabel subsystems and
entanglement without changing any invariant probability.  A relative fact may
be created in an interaction without becoming a durable, manipulable memory.
Neither implication
\begin{equation}
 \text{relational fact}\Rightarrow\text{stored carrier},\qquad
 \text{frame-relative state}\Rightarrow\text{different reachable family}
\end{equation}
holds without an additional physical construction.  A proposed carrier must
either survive the relevant frame equivalences or be explicitly indexed to a
declared frame and observer algebra.

\subsection{The Dirac construction is a rung, not the missing arrow}

The bounded Dirac result belongs between supplied quantum structure and an
effective one-particle role.  Its local update specifications, convergence
tests, causal cones, interventions, and observable witnesses show how a known
role can be recovered constructively \citep{Clark2026DiracRole}.  The hostile
audit and exact time-frame relation limit the architecture-independence claim,
but do not erase the intact shared role core.

Its limitation is decisive for the present paper: complex amplitudes, Hilbert
composition, unitary updates, a Born readout, a physical clock, and a mass
schedule enter as supplied structure.  Thus the construction answers
``can these local quantum updates realize this bounded Dirac role?''  It does
not answer ``can IOTA interactions create quantum state structure?''  A visual
analogy between an IOTA coherence and a wave packet would not close this gap.
The required bridge must recover interference and composition under novel
tests, not encode the target equation in a relabeled update rule.

\subsection{Quasiparticles establish possibility, not provenance}

\paragraph{Established field result.}
Local quantum organization can produce excitations whose operational identity
is not that of any microscopic constituent.  Exactly soluble string-net phases
support emergent gauge and fermionic roles, and the Kitaev honeycomb model
supports fractionalized Majorana and flux excitations
\citep{LevinWen2005StringNet,Kitaev2006Anyons}.  Material measurements provide
model-dependent evidence for such fractional excitations, while controlled
processors and quantum gases have created and manipulated non-Abelian or
anyonized collective behavior
\citep{JansaEtAl2018Fractional,IqbalEtAl2024NonAbelian,
DharEtAl2025Anyonization}.  These are not merely metaphorical particles: within
their domains they have measurable statistics, braiding, correlations, or
response signatures.

One-dimensional quantum cellular automata provide an especially close
existence precedent.  A multiparticle QCA can match its quantum walk in the
one-particle sector and approach a free-fermion QFT at long wavelength.  The
same analysis proves an obstruction for the proposed class in higher
dimensions \citep{MlodinowBrun2020QCAQFT}.  The positive construction and the
no-go result are equally relevant: local updates can recover a field role, but
not generically and not without dimensional and structural conditions.

\paragraph{Synthesis inference.}
These results defeat the blanket objection that microscopic carrier variables
and effective particle identity must coincide.  They make the later portion of
the IOTA ladder logically plausible.  They do not identify the IOTA lattice as
nature's substrate because every cited example starts with a quantum Hilbert
space, Hamiltonian or circuit, selected phase, and ordinary laboratory
resources.

\paragraph{ICP hypothesis and gate.}
The next project object must be more than a long-lived pattern.  A
prospectively declared excitation should have creation and annihilation
operations, finite propagation, dispersion, composition and exchange rules,
symmetry and conserved quantities, scattering or perturbative response, and
robustness under carrier replacement and representation change.  It must then
survive the imported-quantum, lattice-artifact, encoding, and single-target
nulls.  Passing those gates would establish a quasiparticle-like role in the
model.  Only a prior derivation of the quantum structure would make it evidence
for the full substrate claim.

\subsection{Causal order is a boundary on the word ``future''}

The ordinary form of Eq.~\eqref{eq:future-family} presumes a declared temporal
or operational order.  Process-matrix theory and the quantum switch show that
quantum operations can be organized without one fixed global order, and causal
witness experiments have certified causal nonseparability under trusted device
models \citep{OreshkovEtAl2012CausalOrder}.  Circuit order, process
nonseparability, causal-inequality violation, and indefinite order among
localized spacetime events are nevertheless distinct claims
\citep{PaunkovicVojinovic2020CausalOrders,OrmrodEtAl2023SectorialConstraints}.
A 2026 photonic experiment reported a Bell-like causal-order-inequality
violation at 5.7 standard deviations under stated free-intervention and
relativistic-separation assumptions, while retaining locality, fair-sampling,
free-choice, and setup loopholes \citep{GuoEtAl2026CausalOrder}.

\paragraph{Boundary verdict.}
A quantum control can condition the family of later operation sequences.  That
control is not thereby a history-generated record; causal nonseparability is
not thereby indefinite spacetime geometry; and neither result derives
causality from information.  Where order is not fixed, ICP must replace an
ordinary time-indexed future with an explicitly ordered process functional.
The causal-order literature is therefore a formal boundary condition on the
proposal, not evidence for the substrate ontology.

\subsection{Quantum-sector conclusion}

The quantum literature supports three ingredients: interactions can create
physically separated records; record accessibility is relative to an algebra,
factorization, reference frame, and operation class; and organized quantum
substrates can realize genuine particle-like excitations.  The synthesis
inference is that persisted information should be treated as an embodied
source--carrier--operation relation rather than a free-floating substance.
The open IOTA/ICP hypothesis is that such roles can be constructed from the base
informational dynamics and can acquire quantum composition without the quantum
target being supplied.  No result reviewed here closes that arrow.
